\section{现有标定工程与测试思路审查}
\subsection{审查范围与结论分级}
本章直接参考上传的 \path{TiltSensorCalibration_PLC_AutoTempComp_WindowsCompatible(2).zip} 中实际文件，不假设曾在本环境运行 Windows DLL 或控制实体转台。审查以\textbf{源代码与现行试运行配置}为证据；源文件的历史版本说明中“已验证”的含义通常是离线静态/合成测试，\textbf{不是本次硬件现场实测事实}。

\begin{longtable}{C{9mm}p{46mm}p{70mm}p{22mm}}
\toprule
序 & 现状或问题 & 本规范建议 & 判断\\\midrule\endfirsthead
\toprule 序 & 现状或问题 & 本规范建议 & 判断\\\midrule\endhead
1 & 已有 12 ACC 拟合 + 12 独立验证、全矩阵/零偏模型、A/B/C 对比。 & 继续保留；\textbf{新增锁箱验收集}以避免选模型使用验证集后再次宣称独立验收。 & 必改方法\\
2 & \texttt{outer\_x\_inner\_z}，$R=R_x(\alpha)R_z(\beta)$，公共 P 系。 & 符合 U--T 双轴转台的候选运动学思路；但\textbf{仍须实机确认符号/零位}，不可把内轴 Z 误写成第二独立俯仰轴。 & 应核实\\
3 & GYRO 拟合用 $\pm1,\pm2,\pm5^\circ/s$；验证用 $\pm1.5,\pm3.5^\circ/s$。 & 这些点有利于低速诊断，但量程覆盖非常窄；更关键的是手册\textbf{速率精度无要求}，需证书或独立速率基准才能验收绝对比例因子。 & 有条件\\
4 & \texttt{position\_tolerance} $=0.2^\circ$；转台手册给出角位置 $\pm3''$。 & $0.2^\circ$ 是目前\emph{到位/安全状态检查容差}，不可直接作为静态倾角精度参考误差。正式评估必须使用实际位置反馈、夹具外参和误差预算。 & 必改判据\\
5 & 工艺质量筛查：静止 ACC std $\le0.025g$，GYRO std $\le0.45^\circ/s$。 & 这是\emph{可采集性/稳态质量}阈值，未必适用于高精度倾角的噪声要求；应把数据质量门槛和产品测量精度门槛彻底分开。 & 必区分\\
6 & 候选验收：ACC 向量 RMSE $\le0.015g$，方向最大误差 $\le1.2^\circ$。 & 这些值允许的残余可能高于未来高精度倾角目标。未提供产品目标前不得宣称已达到 $0.1^\circ$ 或更高精度。 & 待审批\\
7 & 候选 GYRO 残差最大值 $0.65^\circ/s$，零速残余 $0.15^\circ/s$。 & 作为研发现状记录而非权威标准；必须补充静态零偏、重复性、速率精度可溯源评估。 & 待审批\\
8 & 常温质量配置允许芯片温度 $15$--$35^\circ$C，单工况跨幅 $\le1^\circ$C。 & 当前转台\emph{实验室工作环境}为 $25\pm5^\circ$C，\emph{芯片温度}与\emph{室温}要分开；单工况稳定不等于整轮无温差。 & 必补记录\\
9 & 3$\sigma$/MAD 异常数据筛选、静态最多额外 5 次重试，速率最多额外 3 次。 & 方法适合防止故障值污染拟合，但\textbf{不能只报告通过的最后一轮}；最终验收保留首次失败、所有样本和重试统计，避免结果选择偏差。 & 应改报告\\
10 & 软件允许量程/ODR 未由 MCU 确认时生成 \texttt{PILOT\_PASS\_}\allowbreak\texttt{NOT\_RELEASED}。 & 正确保留试运行而非正式 PASS；执行 S1 时如转换比例未核实，对\emph{绝对数值}测试应标为 \texttt{BLOCKED}。 & 保留保护\\
11 & 16 点帧时间由帧接收时刻/帧间隔估算，转台由线程轮询取主机时间戳。 & 并非硬件同步；动态相位/延迟只能诊断，不能将插值/互相关后误差当作真实系统延迟。 & 必补同步\\
12 & \texttt{alignment} 最新输出使用公共机械平台 P 系补偿，没有再乘一个“相对旋转”。 & 此方向比无条件再做一次轴对齐更合理；对输出需检查坐标统一后是否仍有夹具安装误差、角度定义或重复旋转。 & 保留核实\\
13 & 工程同时包含 2TD600 与 2FDT600WK 的 DLL/温补模块。 & \textbf{本 S1 只按 2TD-600 手册约束。} 不能将温控 2FDT 的位置差分转速估计或温箱联锁能力视为本 2TD-600 已具备。 & 必区分设备\\
14 & 工艺配置 \texttt{approved\_for\_}\allowbreak\texttt{production=false}。 & 文档全部采用“试运行/待审批”结论模式；不把软件产生 JSON 或 Excel 等价于已满足正式规格。 & 保留保护\\
15 & 工艺中的禁用候选速率含 $80^\circ/s$ 验证点；手册两轴速率上限为 $60^\circ/s$。 & \textbf{不得解除候选速率禁用开关后直接执行。} 如要测 $80^\circ/s$，必须另选额定能力足够且已计量的设备，或修改工况。 & 禁止直接使用\\\bottomrule
\end{longtable}
\subsection{建议对现有软件做的最小化调整（非本次源码修改）}
\begin{enumerate}
\item 新增独立运行模式 \texttt{POSTCAL\_VERIFY}：\emph{加载 S0 冻结参数，仅采集和报告}，禁止写入新内参；原有标定模式保持不变。
\item 将三个数据集元信息强制写入报告：\texttt{fit}、\texttt{model\_select}、\texttt{acceptance\_lockbox}，并用 run\_id 与 CRC 防止串组。
\item 评估\texttt{2TD600} 实际速率反馈的时序与计量来源；新增“参考速率未标定”标记，阻止绝对比例因子自动 PASS。
\item 增加双轴实际位置与期望位置分别记录、参考重力矢量/角速度矢量单独存盘、所有失败与重试次数及数据质量分布展示。
\item 输出\texttt{PASS/FAIL/BLOCKED/DIAGNOSTIC} 四态；保留现有\texttt{PILOT\_PASS\_}\allowbreak\texttt{NOT\_RELEASED} 作为工艺审批维度，不与测试结果状态混淆。
\item 将“常温实验室温度”与“芯片内温”作为两个独立通道；将补偿矩阵映射到公共 P 系的约定和未来 PCB--外壳外参分开存储。
\end{enumerate}
\begin{keynote}
\textbf{最重要的修订：} 独立验收不能复用已经参与 A/B/C 模型选择的验证数据；速率精度不能基于厂家没有保证的速率反馈给出正式合格结论；内轴绕 Z 运动在外轴水平时不是第二个独立倾角刺激；常温质量门槛不能替代目标静态角度精度。
\end{keynote}
